\documentclass{article}
\usepackage{Self_Style}

\title{UChicago REU Notes Week 1}
\author{Zih-Yu Hsieh}
\date{\today}

\begin{document}
\maketitle
\tableofcontents

\section{Week 1}
Here it's mainly the preparation for any further theory in spectra. To learn the material, here I focused mainly on some topological preliminaries. Main reference comes from \emph{Concise Course in Algebraic Topology}.

\begin{rem} %changes
I believe most texts assume that compactness means not only the existence of finite subcover for arbitrary open cover, but also Hausdorffness on the space. Here, we also assume this.
\end{rem}

\subsection{Compactly Generated Spaces}

In most literatures I've read, the spaces are assumed to be compactly generated. So, this section is collecting some information I have about it. First, starting with the notion of $k$-spaces:

\begin{defn}
    Given a topological space $X$, a subset $A \subseteq X$ is \emph{$k$-closed}, if for any compact space $K$, and continuous map $g:K\rightarrow X$, the preimage $g^{-1}(A) \subseteq K$ is closed.

    $X$ is a \emph{$k$-space}, if any $k$-closed subset $A \subseteq X$ is closed.
\end{defn}

Which, adapting the usual continuous maps, this realizes a \emph{Category of $k$-spaces}, denoting $k\mathcal{U}$. There are some properties about this category, regarding a functor from regular category of spaces $\mathcal{U}$, and this category's completeness / cocompleteness.

\begin{defn}
    Given a space $X$, define a topological space $kX$ where the underlying set coincides with $X$, while requiring all $k$-closed sets $A \subseteq X$ to be closed in $kX$ as its new topology.
\end{defn}

This construction satisfies the closed sets axioms, hence form a topology; also, under the new topology, $kX$ is a $k$-space. Furthermore, since all closed sets are automaticlaly $k$-closed, the set identity function $\iota_X:kX\rightarrow X$ is in fact continuous.

We want to further argue $k$ is a functor $\mathcal{U}\rightarrow k\mathcal{U}$, which can be achieved with the following proposition:

\begin{prop}
Let $Z$ be a $k$-space, $Y$ arbitrary topological space, and $f:Z\rightarrow Y$ a continuous map. Then, there exists a unique continuous map $kf:Z\rightarrow kY$, such that the following diagram commutes:
\[\xymatrix{
    Z \ar@{-->}[dr]_-{\exists ! kf} \ar[rr]^-{\forall f} && Y\\
    & kY \ar[ur]_-{\iota_Y}
} \]
\end{prop}
\begin{proof}
    Since $\iota_Y$ is a set identity map, uniqueness of $kf$ is clear (as set function it coincides with $f$). To show $kf$ is continuous, it suffices to show all $k$-closed set $A\subseteq Y$ has preimage $f^{-1}(A)\subseteq Z$ being closed.

    Let $g:K\rightarrow Z$ be arbitrary continuous map, with $K$ being compact. Consider $f\circ g:K\rightarrow Z\rightarrow Y$, for any $A\subseteq Y$ that's $k$-closed, $(f\circ g)^{-1}(A)= g^{-1}(f^{-1}(A))\subseteq K$ is closed, hence $f^{-1}(A)\subseteq Z$ is also $k$-closed. Since $Z$ is a $k$-space, $f^{-1}(A)$ is closed.
\end{proof}

\begin{cor}
Given any continuous map $f:X\rightarrow Y$, there exists a unique continuous map $kf:kX\rightarrow kY$, such that the following diagram commutes:
\[\xymatrix{
    X \ar[r]^-{\forall f} & Y\\
    kX \ar[u]^-{\iota_X} \ar@{-->}[r]_-{\exists ! kf} & kY \ar[u]_-{\iota_Y}
} \]
\end{cor}
\begin{proof}
    Since $f\circ \iota_X:kX\rightarrow Y$ is continuous, with $kX$ being a $k$-space, \textbf{Proposition 1.4} guarantees the existence and uniqueness of $kf$ (which as set map it also agrees with $f$, as $\iota_X$ is set identity).
\end{proof}
This formulates $k:\mathcal{U}\rightarrow k\mathcal{U}$ as a functor. And, if consider the inclusion functor $\mathcal{I}:k\mathcal{U}\rightarrow\mathcal{U}$, \textbf{Proposition 1.4} also insists $(\mathcal{I},k)$ as an adjoint pair, as it phrases a natural isomorphism as follow:
\[\Hom_\mathcal{U}(\mathcal{I}Z, Y)\cong \Hom_{k\mathcal{U}}(Z,kY),\quad f\mapsto kf\]
Using some categorical nonsense, $k$ as a right adjoint preserves all limits, so $k\mathcal{U}$ is a complete category. In paticular, we denote $X\times^k Y:=k(X\times Y)$ as the product of two $k$-spaces $X,Y$. 

To show it's also cocomplete, consider the following:
\begin{prop}
    Quotient of a $k$-space is still a $k$-space; arbitrary disjoint union of $k$-spaces is a $k$-space.
\end{prop}
\begin{proof}
    Let $X$ be a $k$-space, and $f:X\rightarrow Y$ be a quotient map. For any $k$-closed subset $A\subseteq Y$, consider arbitrary continuous map $g:K\rightarrow X$ with $K$ compact. Given the composition $f\circ g:K\rightarrow X\rightarrow Y$, $(f\circ g)^{-1}(A) = g^{-1}(f^{-1}(A))\subseteq K$ is closed, hence $f^{-1}(A)\subseteq X$ is $k$-closed, which is closed as $X$ is a $k$-space. Then, using the quotient space property, $A = f(f^{-1}(A))\subseteq Y$ is closed, concluding that $Y$ is also a $k$-space.

    \hfil

    Now, given $\{X_i\}_i$ a family of $k$-spaces, consider $Z:= \coprod_i X_i$, with $\iota_i:X_i\rightarrow Z$ the canonical inclusion. Given any $k$-closed subset $A\subseteq Z$, similar logic from previous proofs show that $A\cap X_i = \iota_i^{-1}(A)\subseteq X_i$ are all $k$-closed, hence closed (as $X_i$'s are $k$-spaces). Which, this verifies $A\subseteq Z$ is closed under the disjoint union topology, showing $Z$ is also a $k$-space.
\end{proof}
This guarantees the corproduct and coequalizers of $k$-spaces (and maps between them) in $\mathcal{U}$ are still $k$-spaces, which these colimits exist in $k\mathcal{U}$, showing the cocompleteness.

\hfil

Now, move on to another criterion for compactly generated spaces:
\begin{defn}
    A space $X$ is \emph{$k$-Hausdorff} if the diagonal $\Delta X\subset k(X\times X)$ is closed.
\end{defn}

Denote $h\mathcal{U}$ as the \emph{Category of $k$-Hausdorff spaces}. 

There's a lot more details for this, I'll try and fill it in afterward==.

There is another functor $h:\mathcal{U}\rightarrow h\mathcal{U}$ which generates a $k$-Hausdorff space out of any topological spaces, and with inclusion functor $\mathcal{J}:h\mathcal{U}\rightarrow \mathcal{U}$, they form an adjoint pair $(h,\mathcal{J})$. So, $h$ as a left adjoint preserves colimits. 

\hfil

Finally, for compactly generated spaces, it's defined as a $k$-Hausdorff $k$-space. In this space, limits and colimits exist, but need to be constructed in $\mathcal{U}$, then apply the functor $h$ or $k$ to yield the desired results.

Another equivalent way to phrase $k$-Hausdorffness when given a $k$-space, is the \emph{weak Hausdorffness}: Given any continuous map $g:K\rightarrow X$ with $K$ compact, $g(K)\subseteq X$ is a closed subset.

\hfil

\subsection{Spaces of Maps}
This section, we'd like to realize $Y^X:=\Hom_\Top(X,Y)$ (the set of all continuous maps $X\rightarrow Y$) as a topological space. Here, we consider the \emph{compact open topology}.

First, given the any map $f:K\times X\rightarrow Y$, there is a corresponding set function $\tilde{f}:K\rightarrow Y^X$ by $\tilde{f}(k) = f(k,\_):X\rightarrow Y$. 
\begin{defn}
    The \emph{compact open topology} on $\Hom_\Top (X,Y)$, has $S\subseteq Y^X$ being open iff all compact space $K$, continuous map $f:K\times X\rightarrow Y$, the set 
    \[\tilde{f}^{-1}(S) = \{k\in K | f(k,\_)\in S\}\]
    is open in $K$.
\end{defn}
Obviously, $\emptyset, Y^X$ are both open; if $\{S_i\subseteq Y^X\}_i$ are open, then $\tilde{f}^{-1}(\bigcup_i S_i)=\bigcup_i \tilde{f}^{-1}(S_i)$ is open in $K$, so $\bigcup_i S_i$ is open; finally, if $S,S'\subseteq Y^x$ are open, then $\tilde{f}^{-1}(S\cap S')=\tilde{f}^{-1}(S)\cap \tilde{f}^{-1}(S')$ is open in $K$, so $S\cap S'$ is also open. This defines a topology.

Here is some nice properties of composition of maps:
\begin{prop}
    Given $g:X'\rightarrow X$, the precomposition $g^*:=\_\circ g:Y^X\rightarrow Y^{X'}$ is continuous. Dually, given $h:Y\rightarrow Y'$, the postcomposition $h_*:=h\circ\_:Y^X\rightarrow Y'^X$ is continuous.
\end{prop}
\begin{proof}
    In the first scenario, given arbitrary open subset $S'\subseteq Y^{X'}$, the preimage $S:=(g^*)^{-1}(S') = \{\ell\in Y^X|\ell\circ g\in S'\}$. Fix arbitrary continuous map $f:K\times X\rightarrow Y$ with $K$ compact, consider $\tilde{f}^{-1}(S)\subseteq K$, we get:
    \[\tilde{f}^{-1}(S) = \{k\in K|f(k,\_)\in S\} = \{k\in K|f(k,\_)\circ g\in S'\}\]
    Similarly, consider $f' := f\circ (\id_K\times g):K\times X'\rightarrow Y$ (by $f'(k,x') = f(k,g(x'))$), the set $\tilde{f'}^{-1}(S')\subseteq K$ is open by assumption, which it has the following description:
    \[\tilde{f'}^{-1}(S')=\{k\in K|f'(k,\_)\in S'\} = \{k\in K|f(k,\_)\circ g\in S'\} = \tilde{f}^{-1}(S)\] 
    This shows the openness of $S\subseteq Y^X$.

    \hfil

    For the second scenario, given arbitrary open subset $S'\subseteq Y'^X$, the preimage $S:= h_*^{-1}(S') = \{\ell\in Y^X|h\circ \ell\in S'\}$. Fix arbitrary continuous map $f:K\times X\rightarrow Y$ with $K$ compact, notice the following description:
    \[\tilde{f}^{-1}(S) = \{k\in K|f(k,\_)\in S\} = \{k\in K|h\circ f(k,\_)\in S'\}\]
    Which, notice $h\circ f:K\times X\rightarrow Y\rightarrow Y'$ satisfies $\tilde{(h\circ f)}^{-1}(S') = \{k\in K|h\circ f(k,\_)\in S'\} = \tilde{f}^{-1}(S)$, and by assumption this set is open. So, $S\subseteq Y^X$ is open, finishing the claim.
\end{proof}

As a result, the Hom functor $\Hom_\Top:\Top^\Op \times \Top\rightarrow \Top$ is a bifunctor.

\hfil

If restrict to the compactly generated spaces, there is a homeomorphism \[\Hom_\Top (X\times^k Y, Z)\cong \Hom_\Top (X, \Hom_\Top (Y,Z))\]
Yeah, more detailes to fill in.

\subsection{Cofibration}

Intuitively, these are the maps where homotopy can be extended to the codomain:
\begin{defn}
    A map $\iota:A\rightarrow X$ is a \emph{cofibration}, if it satisfies the \emph{Homotopy ExtensionProperty} (HEP). This means, given any map $f:X\rightarrow Y$, homotopy $H:A\times I\rightarrow Y$, zeroth inclusions $i_0:A\rightarrow A\times I$ and $i_0:X\rightarrow X\times I$ (by $x\mapsto (x,0)$), and product map $\iota\times \id_I:A\times I\rightarrow X\times I$, the homotopy extends to $\tilde{H}:X\times I\rightarrow Y$, with the following commutes:
    \[\xymatrix{
    A \ar[dd]_-{\iota} \ar[rr]^-{i_0} && A\times I \ar[dl]_-{H} \ar[dd]^-{\iota\times \id_I}\\
    & Y\\
    X \ar[ur]^-{f} \ar[rr]_-{i_0} && X\times I \ar@{-->}[ul]_-{\exists \tilde{H}}
} \]
\end{defn}
Intuitively, it gives a description of how to map $X$ into $Y$ (also how to map $A$ into $Y$), generates a homotopy $H$ on $A$ only, and the property guarantees it to extend to a homotopy on $X$.

It's preserved under fibre product:
\begin{lem}
    If $\iota:A\rightarrow X$ is a cofibration, $g:A\rightarrow B$ any map, then the induced pushout $\iota':B\rightarrow B\sqcup_A X$ is a cofibration:
    \[\xymatrix{
        A \ar[d]_-{g} \ar[r]^-{\iota} & X \ar[d]^-{g'}\\
        B \ar[r]_-{\iota'} & B\sqcup_A X
    }\]
\end{lem}
\begin{proof}
    Before starting, notice that product preserves adjunction/fibre coproduct, because $(B\sqcup_A X)\times I\cong (B\times I)\sqcup_(A\times I) (X\times I)$ (i.e. the first space can be viewed as the disjoint union of the product, then glue using identical methods at each $t\in I$).

    Then, given map $f:B\sqcup_A X\rightarrow Y$, homotopy $H:B\times I$, and the fibre coproduct maps, we get:
    \[\xymatrix{
        A \ar[dr]^-{g} \ar[dddd]_-{\iota} \ar[rrrr]^-{i_0} &&&& A\times I \ar[dddd]^-{\iota\times \id_I} \ar[dl]_-{g\times \id_I}\\
        & B \ar[dd]_-{\iota'} \ar[rr]^-{i_0} && B\times I \ar[dd]^-{\iota'\times \id_I} \ar[dl]_-{H}\\
        && Y\\
        & B\sqcup_A X \ar[ur]^-{f} \ar[rr]_-{i_0} && (B\sqcup_A X)\times I\\
        X \ar[ur]^-{g'} \ar[rrrr]_-{i_0} &&&& X\times I \ar[ul]_-{g'\times \id_I} \ar@{-->}@(l,d)[uull]^-{\exists \tilde{H}}
    }\]
    The existence of $\tilde{H}$ is given by HEP of $\iota:A\rightarrow X$. This induces the following diagram:
    \[\xymatrix{
        A\times I \ar[r]^-{\iota\times \id_I} \ar[d]_-{g\times \id_I} & X\times I \ar[d]^-{g'\times \id_I} \ar@(r,u)[ddr]^-{\tilde{H}}\\
        B\times I \ar@(d,l)[drr]_-{H} \ar[r]_-{\iota'\times \id_I} & (B\sqcup_A X)\times I \ar@{-->}[dr]^-{\exists ! \overline{H}}\\
        && Y
    }\]
    where $\overline{H}$ is induced by the fibre coproduct property. Plugin to the previous diagram, we yield the desired homotopy extension for the inner square.
\end{proof}

\hfil

If consider the ``special case'' $Y= M_\iota$, the mapping cylinder of $\iota:A\rightarrow X$ (which is constructed as $M_\iota := (A\times I\sqcup X)/\sim$, where $(a,0)\sim f(a)$). Notice it generates the following fibre coproduct diagram:
\[\xymatrix{
    A \ar[d]_-{\iota} \ar[r]^-{i_0} & A\times I \ar[d]^-{H}\\
    X \ar[r]_-{i_0'} & M_\iota
}\]
So, the HEP and fibre coproduct property guarantees the following diagram:
\[\xymatrix{
    A \ar[dd]_-{\iota} \ar[rr]^-{i_0} && A\times I \ar[dl]_-{H} \ar[dd]^-{\iota\times \id_I}\\
    & M_\iota \ar@{-->}@/_{.5pc}/[dr]_-{\exists ! j}\\
    X \ar[ur]^-{i_0'} \ar[rr]_-{i_0} && X\times I \ar@{-->}@/_{.5pc}/[ul]_-{\exists \tilde{H}}
} \]
Here, $\tilde{H}$ is the homotopy extension, $j$ is the factorization by fibre coproduct. And, this enforces $\tilde{H}\circ j = \id_{M_\iota}$, showing $j$ is injective. If restrict to $t=1$, then $j:A\times \{1\}\rightarrow X\times\{1\}$ is precisely the map $\iota$, showing it's injective (in fact it's an embedding, because we also have a retract $\tilde{H}:X\times\{1\}\rightarrow A\times\{1\}$).

From other sources I can find, $\iota$ unfortunately doesn't have closed image in general, yet for the case of $A,X$ being compactly generated (for instance, CW Complexes) I believe this is true. (Maybe by $k$-closed property).

\hfil

So, we see $\iota:A\rightarrow X$ is a cofibration (for compactly generated spaces) directly implies $\iota$ is an embedding with closed image (or $A$ is a closed subspace of $X$); and, the mapping cylinder is precisely $M_\iota=X\times \{0\}\cup A\times I$ in $X\times I$, the extended homotopy guarantees it's a retract of $X\times I$.

Conversely, given $A\subseteq X$ a closed subspace, while its mapping cylinder $M_\iota\subseteq X\times I$ is a retract of $X\times I$. Then, use the retract as the extended homotopy (together with the fact that it's a fibre coproduct, hence the ``initial extension diagram'' whenever seeking the homotopy extension for other $Y$), we get $A\hookrightarrow X$ is a cofibration.

So, for compactly generated spaces, it's equivalent to say: a cofibration is a closed embedding $A\subseteq X$, such that the mapping cylinder $X\times \{0\}\cup A\times I$ is a retract of $X\times I$.

\hfil

To talk about homotopy equivalences of cofibers (which are not necessarily cofibrations, but just maps $A\rightarrow X$ when fixing $A$). Given, $i:A\rightarrow X$, $j:A\rightarrow Y$ two cofibers, a map $f:X\rightarrow Y$ is a morphism, if the following diagram commutes:
\[\xymatrix{
    & A \ar[dl]_-{i} \ar[dr]^-{j}\\
    X \ar[rr]_-{f} && Y
}\]

\begin{defn}
    Two such morphisms $f,f':X\rightarrow Y$ are homotopic in this category (under $A$), if they're homotopic relative to $A$, i.e. there exists homotopy $H:X\times I\rightarrow Y$ of $f,f'$, where all $a\in A, t\in I$ satisfies $H(i(a),t)=j(a)$ (namely $A$ is fixed under this homotopy). This is called \emph{cofiber homotopy}.

    The homotopy equivalence built on cofiber homotopies is called \emph{cofiber homotopy equivalence}.
\end{defn}

This notion of homotopy equivalence is much stronger than original homotopy. Yet, when restricting to cofibrations, then it doesn't make a difference:
\begin{prop}
    Given $i:A\rightarrow X$, $j:A\rightarrow Y$ two cofibrations, if morphism $f:X\rightarrow Y$ is a homotopy equivalence, then it's a cofiber homotopy equivalence.
\end{prop}

\begin{proof}
    It suffices to find $g:Y\rightarrow X$ that's a homotopy invese, with $g\circ f\simeq \id_X$ relative to $A$ (because then apply the same argument on $g:Y \rightarrow X$, there is a homotopy inverse $f':X\rightarrow Y$, such that $f'\circ g\simeq \id_Y$ relative to $A$. As a result, $f = \id_Y \circ f\simeq f'\circ g\circ f\simeq f'\circ \id_X = f'$ relative to $A$, hence $f\circ g\simeq f'\circ g\simeq id_Y$ relative to $A$, showing it's a cofiber homotopy equivalence).

    First, by hypothesis, there is $g'':Y\rightarrow X$ a homotopy inverse of $f$. (Not done)
\end{proof}

\hfil

\subsection{Fibration}

As a dual notion to cofibration (which ``extends'' homotopy to the target space), fibration is the one that ``lifts'' the homotopy to the source space. 

\begin{defn}
    A surjective map $p:E\twoheadrightarrow B$ is a \emph{fibration}, if it satisfies \emph{covering homotopy property} (CHP in May's book), or \emph{homotopy lifting property}. Meaning, given any $f:Y\rightarrow E$, and homotopy $H:Y\times I\rightarrow B$ of $p\circ f:Y\rightarrow E\twoheadrightarrow B$, there exists a homotopy $\tilde{H}:Y\times I\rightarrow E$ of $f$, with the following diagram comutes:
    \[\xymatrix{
        Y \ar[d]_-{i_0} \ar[r]^-{\forall f} & E \ar@{->>}[d]^-{p}\\
        Y\times I \ar@{-->}[ur]^-{\exists \tilde{H}} \ar[r]_-{\forall H} & B
    }\]
\end{defn}
A typical example is the covering map, where homotopies can be lifted uniquely once a basepoint of the covering is chosen. So, fibration can be viewed as a generalization of covering space, where homotopy liftings exist.

Using the $k$-ified compact-open topology of the function space (in particular, the homeomorphism $\Hom_\Top (K\times X,Y)\cong \Hom_\Top (K,\Hom_\Top (X,Y))$, and $Y^X:=\Hom_\Top (X,Y)$), the above can be reformulated as the following (which is structure wise more similar to the dual of the cofibration diagram):
\[\xymatrix{
    E \ar@{->>}[dd]_-{p} && E^I \ar[ll]_-{p_0} \ar[dd]_-{p\circ \_}\\
    & Y \ar[ul]^-{\forall f} \ar[dr]_-{\forall H} \ar@{-->}[ur]^-{\exists \tilde{H}}\\
    B && B^I \ar[ll]^-{p_0}
}\]
Then, by dualizing the proof for cofibration, we yield the following result:
\begin{lem}
    Pullback of fibration is fibration: Given fibration $p:E\twoheadrightarrow B$, and any map $g:B'\rightarrow B$, the base change $p':E\times_B B' \rightarrow B'$ is a fibration:
    \[\xymatrix{
        E\times_B B' \ar[r]^-{g'} \ar@{->>}[d]_-{p'} & E \ar@{->>}[d]^-{p}\\
        B' \ar[r]_-{g}& B
    }\]
\end{lem}

\hfil

Similar to the existence of mapping cylinder $M_\iota$ as the ``initial test diagram'' for cofibration (as it's the fibre coproduct of cofibration $A\rightarrow X$, and inclusion $A\rightarrow A\times I$), it suggests the exitence of ``final test diagram'' for fibration, when choosing $Y=N_p := E\times_B B^I$ (the fibre product of fibration $E\rightarrow B$, and evaluation by zero map $B^I\rightarrow B$). The fibre product property (together with the test diagram of fibration) guarantees $N_p$ is the final diagram (namely, all other choices of $Y$ factors uniquely through $N_p$).

A precise description of $N_p$ is as follow:
\[N_p = \{(e,\beta) | \beta(0)=p(e)\} \subseteq E\times B^I\]  

Plugin to the diagram, we observe the following using the fibration and fibre product property:
\[\xymatrix{
    E \ar@{->>}[dd]_-{p} && E^I \ar@{-->}@/_{.5pc}/[dl]_-{\exists ! k} \ar[ll]_-{p_0} \ar[dd]^-{p\circ \_}\\
    & N_p \ar[ul]^-{p_1} \ar[dr]_-{p_2} \ar@{-->}@/_{.5pc}/[ur]_-{\exists s}\\
    B && B^I \ar[ll]^-{p_0}
}\]
Here, $k$ is induced by fibre coproduct, $s$ is induced by fibration property. Universality of $N_p$ (as fibre product) guarantees that $k\circ s = id_{N_p}$, another name for these lifting $s$ is called \emph{path lifting function} (as it lifts a pair $(e,\beta)\in N_p$ with $\beta(0)=p(e)$, to a path $\tilde{beta}\in E^I$, that satisfies $\tilde{\beta}(0)=e$, and $p\circ \tilde{beta}=\beta\in B^I$).

A special case - when $p:E\rightarrow B$ is a covering, then $s$ is in fact unique, due to the unique path lifting property (whenever fixing a path and a basepoint). 

\hfil

For common types of fibrations, the case for covering space suggests the notion of fibre bundle could work: Given covering $p:E\rightarrow B$, every point $x\in B$ has an open neighborhood $U_x\subseteq B$, such that $p^{-1}(U_x)\cong U_x\times p^{-1}(x)$ (where $p^{-1}(x)$ is a discrete subspace), and all fibres are homeomorphic to each other (i.e. same cardinality when the fibre is discrete).

If restricts the bundle to have \emph{numerable open covers} (namely, an open cover $\{U_i\}_i$, such that each point $x\in B$ is contained in finitely many $U_i$'s, and each $U_i$ is equipped with a continuous function $\lambda_i:B\rightarrow I$ such that $\lambda_i^{-1}(0,1] = U_i$). Then, the following theorem holds:
\begin{thm}
Given a map $p:E\rightarrow B$ and $\{U_i\}_i$ a numerable open cover of $B$. Then, $p$ is a fibration iff the restrictions $p|_{U_i}:p^{-1}(U_i)\rightarrow U_i$ is a fibration for every $U_i$.
\end{thm}

This guarantees all fibre bundles with numerable trivialization cover to be a fibration.

\hfil

Similar to cofibre homotopy, we define fibre homotopy as follow:
\begin{defn}
    Given two fibres $p:D\rightarrow B$ and $p':E\rightarrow B$ (Note: they're not necessarily fibrations), a map $f:D\rightarrow E$ is a morphism if the following diagram commutes:
    \[\xymatrix{
        D \ar[dr]_-{p} \ar[rr]^-{f} && E \ar[dl]^-{p'}\\
        & B
    }\]
    And, given two morphisms $f,f':D\rightarrow E$, they're \emph{fibre homotopic} if there is a homotopy $H:D\times I\rightarrow E$, such that $p = p'\circ H:D\rightarrow B$ (namely, at all time $t\in I$, the map preserves $p$).
\end{defn}

By dualizing the notion fro cofibre, we again have the following criterion of fibre homotopy equivalence:
\begin{prop}
    Let $p:D\rightarrow B$, $q:E\rightarrow B$ be two fibrations. If $f:D\rightarrow E$ be a homotopy equivalence morphism (i.e. $q\circ f=p$), then it's a fibre homotopy equivalence.
\end{prop}

\hfil

To talk about the relations between the fibres (for instance, for fibre bundle all fibres are homeomorphic), the use of paths (as homotopy between points) can also b lifted given fibration, and can be used to study such relation.

Let $p:E\rightarrow B$ be a $\beta:I\rightarrow B$ be a path with $\beta(0)=b$ and $\beta(1)=b'$, and $F_x:=p^{-1}(x)\subseteq E$ as the fibre (as subspace under $E$). Then, one can consider the following diagram, using the lifting property:
\[\xymatrix{
    F_b\times \{0\}\cong F_b \ar[dd]_-{i_0} \ar[rr]^-{\iota_b} && E \ar[dd]^-{p}\\
    \\
    F_b\times I \ar@{-->}[uurr]^-{\exists \tilde{\beta}} \ar[r]_-{\pi_2} & I \ar[r]_-{\beta} & B
}\]
Which, specify to $t=1$ for $\tilde{\beta}$, we realize the following:
\[p\circ \tilde{\beta}_1 (F_b) = \beta\circ \pi_2 (F_b\times\{1\}) = \beta(1)=b'\implies \tilde{\beta}_1(F_b)\subseteq F_{b'}\]
So, it induces a map $\tilde{\beta}_1:F_b\rightarrow F_{b'}$ (more generally, it induces a continuous family of maps $\tilde{\beta}_t:F_b\rightarrow F_{\beta(t)}$.

\begin{prop}
    Suppose paths $\beta,\beta':I\rightarrow B$ are homotopic, then the induced maps $\tilde{\beta}_1,\tilde{\beta'}_1:F_b\rightarrow F_{b'}$ are also homotopic.
\end{prop}
\begin{proof}
    Let $H:I\times I\rightarrow B$ be the path homotopy of $\beta\simeq \beta'$, so that $H(s,0)=\beta(s)$, $H(s,1)=\beta'(s)$, $H(0,t)=b$, and $H(1,t)=b'$. If consider the subspace of $I\times I$, say $J := (I\times \{0\})\cup (I\times \{1\})\cup (\{0\}\times I)$ (in the homotopy $H$, first side represents $\beta$, second side represents $\beta'$, and the third side represents the side that's constant as $b$), notice there is an automorphism on $I\times I$, so that $J$ gets mapped to $I\times \{0\}$, so the lifting property can be used whenever a map is defined on $J\subset I\times I$.

    Now, define the map map $f:F_b\times J\rightarrow E$, so that on $I\times \{0\}$ it's $\tilde{\beta}:F_b\times I\cong F_b\times I\times \{0\}\rightarrow E$, on $I\times \{1\}$ it's $\tilde{\beta'}:F_b\times I\cong F_b\times I\times \{1\}\rightarrow E$, and on $\{0\}\times I$ it's the map $\iota_b\circ p_1:F_b\times I \cong F_b\times \{0\}\times I \rightarrow F_b \rightarrow E$ (more precisely, we only consider the value in $F_b$, and include it into $E$). By the defining property of each $\tilde{\beta},\tilde{\beta'}$, these three maps glue together, hence $f$ is well-defined. 

    Then, consider the inclusion $i:J\hookrightarrow I\times I$, we form the following diagram using the lifting property:
    \[\xymatrix{
        F_b\times J \ar[dd]_-{\id_{F_b}\times i} \ar[rr]^-{f} && E \ar[dd]^-{p}\\
        \\
        F_b\times I\times I \ar@{-->}[uurr]^-{\exists \tilde{H}} \ar[r]_-{\pi_{2,3}} & I\times I \ar[r]_-{H} & B
    }\]
    Which, if restrict to $F_b\times I\times \{0\}$, $\tilde{H}$ restricts to $\tilde{\beta}$ by our chohice of $f$; similarly, restricting to $F_b\times I\times \{1\}$, $\tilde{H}$ restricts to $\tilde{\beta'}$, showing the two homotopic. So, specify to $t=1$, we get $\tilde{\beta}_1\simeq \tilde{\beta'}_1$, showing the map $[\tilde{\beta}_1]:F_b\rightarrow F_{b'}$ is a well-defined map between fibres, and dependent only on homotopy classes of paths $[\beta]:I\rightarrow B$, connectinb $b$ to $b'$.
\end{proof}

\begin{cor}
    Given a fibration, any two points in the same path component have homotopy equivalent fibres.
\end{cor}
\begin{proof}
    This is because if $\beta,\gamma:I\rightarrow B$ are two paths such that $\beta(1)=\gamma(0)$, then if $b:= \beta(0), b':= \beta(1)=\gamma(0), b'':=\gamma(1)$, then $\tilde{\beta}_1:F_b\rightarrow F_{b'}$ and $\tilde{\gamma}:F_{b'}\rightarrow F_{b''}$ composes to be the same as $\tilde{(\gamma\cdot \beta)}_1:F_b\rightarrow F_{b''}$ (it's induced by the concatination of maps). So, $[\tilde{\gamma\cdot \beta}_1] = [\tilde{\gamma}_1]\circ [\tilde{\beta}_1]$ when passing to homotopy classes of maps.

    Then, since a constant path $c_b$ incudes identity $[\tilde{c_b}_1]:F_b\rightarrow F_b$, we see $[\tilde{\beta^{-1}}_1]\circ [\tilde{\beta}_1]=[\tilde{\beta^{-1}\dot \beta}_1] = [\tilde{c_p}] = [id_{F_b}]$, showing $[\tilde{\beta}_1]$ is a homotopy equivalence.
\end{proof}

\subsection{Based Homotopy}

\begin{defn}
Given a based topological space $X$, its \emph{reduced/based suspension} $\Sigma X:= X\times I/X \vee I$, where $X\vee I$ is the wedge sum of $X$ and inteveral $I$ once fixed a basepoint $x_0\in X$ and $i\in I$, while $X\times I$ has basepoint $(x_0,i)\in X\times I$.

Equivalently, $\Sigma X \cong X\wedge S^1:= X\times S^1/X\vee S^1$.
\end{defn}


\end{document}